\documentclass[11pt]{article}

% \providecommand{\answerDirectory}{solutions}

% Local override, kept so this homework's PDF matches its original exactly:
% the shared macros (src/macros/macros.tex) now define \REPLACEME differently (without the leading space, plus \xspace).
% Defined before ../../macros/macros.tex with \providecommand so a gradescope wrapper's version still wins.
\providecommand{\REPLACEME}{ \textcolor{red}{\ensuremath{\star \cdots \star}}}

../../macros/macros.tex


\inputAnswer{problem_0_answer}

\doHwHeader{Homework \BLUE{8}}

%%%%% EDIT THE FILES NAMED problem_X_answer.tex.  DON'T EDIT THIS FILE

%%%%%%%%%%%%%%%%%%%%%%%%%%%%%%%%%%%%%%%%%%%%%%%%%%%%%%%%%%%% 
%%%%%%%%%%%%%%%%%%%%%%%%%%%%%%%%%%%%%%%%%%%%%%%%%%%%%%%%%%%% 
%%%%%%%%%%%%%%%%%%%%%%%%%%%%%%%%%%%%%%%%%%%%%%%%%%%%%%%%%%%% 
%%%%%%%%%%%%%%%%%%%%%%%%%%%%%%%%%%%%%%%%%%%%%%%%%%%%%%%%%%%% 

\begin{document}

\gradescopeInit


%%%%% EDIT THE FILES NAMED problem_X_answer.tex.  DON'T EDIT THIS FILE 

\paragraph{General instructions.}
% First read \BLUE{Lecture Note 5}, then
Answer the problems given in the rest of this document.
Edit this LaTeX template to add your answers, as described in the file named \texttt{00\_README.txt},
then compile it to produce \BLUE{\texttt{homework\_8.pdf}}\ifGradescope{} for submission to gradescope\fi.

\ifGradescope
For this homework, when you submit to gradescope you don't need to identify the pages
for each problem.  We're using gradescope's ``fixed-length'' grading mode.
To make it work, your answer for each part should fit in the allotted space,
so that each subsequent answer is on the expected page in the expected location.
Your PDF should have
\BLUE{14}
pages.  
\fi

\paragraph{Don't know?  Say so.}
If you don't know the answer to a problem or any part of the problem,
we encourage you to simply write \emph{``I don't know''} instead of giving an answer.
(If you write ``I don't know'' you can also add any questions you might have.)

\paragraph{Gentle reminder: homework is for learning.}
As noted in \LN{prelim}, engaging fully with the homeworks is the best way to learn the material.
\emph{The primary role of homeworks is to help you learn the material,
  and the primary role of feedback on homeworks is to help you improve your ideas.}

Please try to explain your ideas as clearly as you can.
The clearer your explanations are, the more useful the feedback on them can be.

\paragraph{Collaboration and citation policy.}
If you are doing this homework as part of a course, follow the course's collaboration and citation policy.
In any case, at the end of each answer, list the people you collaborated with,
and cite every source that your answer uses ideas from
(other than the lecture notes and the sources listed in the lecture notes).
If there are none, write ``none''.

%%%%% EDIT THE FILES NAMED problem_X_answer.tex.  DON'T EDIT THIS FILE 

\vfill
\noindent{\footnotesize\copyrightNotice}

%%% Local Variables:
%%% mode: latex
%%% TeX-master: "homework_2"
%%% End:


%%%%%%%%%%%%%%%%%%%%%%%%%%%%%%%%%%%%%%%%%%%%%%%%%%%%%%%%%%%% 
%%%%%%%%%%%%%%%%%%%%%%%%%%%%%%%%%%%%%%%%%%%%%%%%%%%%%%%%%%%% 
%%%%%%%%%%%%%%%%%%%%%%%%%%%%%%%%%%%%%%%%%%%%%%%%%%%%%%%%%%%% 
%%%%%%%%%%%%%%%%%%%%%%%%%%%%%%%%%%%%%%%%%%%%%%%%%%%%%%%%%%%% 

%%%%% EDIT THE FILES NAMED problem_X_answer.tex.  DON'T EDIT THIS FILE 

%%%%%%%%%%%%%%%% PROBLEM 1  %%%%%%%%%%%%%%%%%

\newpage 

\problem\emph{(Lecture Note 7)}
You are a freelancer who has to choose which jobs to take on each of $n$ days.  
On each day $i\in\{1,2,\ldots,n\}$ you have three options---take a \emph{light} job, take a \emph{heavy} job, or take no job.
Your pay for the day $i$ is $\ell_i$ if you take a light job, $h_i$ if you take a heavy job, and nothing if you take no job.
There is one additional constraint: for each day $i\in\{2, 3, \ldots, n\}$,
you can take a heavy job on day $i$ only if you took no job in day $i-1$.
(You can take a heavy job on day 1.)

\smallskip

Fix an input $(\ell=(\ell_1,\ldots,\ell_n), h=(h_1,\ldots,h_n))$.

For each $j\in\{0,1,2,\ldots,n\}$, 
a \emph{schedule} for days $\{1, 2, \ldots, j\}$
selects, for each of those $j$ days, one of the three job options (light, heavy, none).
The schedule is \emph{feasible} if, for every day $i\in\{2, 3, \ldots, j\}$,
if the schedule selects the heavy job for day $i$,
then it selects no job for day $i-1$.
The \emph{earnings} of the schedule is the sum of its earnings for each day $i\in\{1,2, \ldots, j\}$,
where the earnings for day $i$ is $\ell_i$ if the schedule selects a light job that day,
$h_i$ if the schedule selects a heavy job, and nothing if it selects no job.
(You may assume $\ell_i$ and $h_i$ are non-negative, but \emph{not} that $\ell_i \le h_i$.)
Two schedules are \emph{distinct} if they are not identical
(i.e., for at least one day, their selections differ).

For $j\in\{0,1,2,\ldots,n\}$, 
define $N(j)$ to be the number of distinct feasible schedules for days $\{1,\ldots, j\}$.
Define $M(j)$ to be the maximum, over all feasible schedule for days $\{1,\dots,j\}$,
of the total earnings of the schedule.


\begin{enumerate}[label=(\alph*)]
\item

Give a recurrence relation for $N(j)$ in terms of $N(0), N(1), \ldots, N(j-1)$.

\item
  
Explain precisely in at most four lines why your recurrence relation is correct. 

\item

Give a recurrence relation for $M(j)$ in terms of $M(0), M(1), \ldots, M(j-1)$, $\ell$, and $h$.

\item
  
Explain precisely in at most four lines why your recurrence relation is correct. 
 
\item

  What is the asymptotic worst-case running time of the resulting (iterative, or recursive and memoized) algorithm for computing $M(n)$, given $(\ell, h)$?

\item
  
  Explain precisely in at most three lines why the bound you give is correct. 

\item

Give pseudo-code for the resulting iterative dynamic-programming algorithm, for computing $M(n)$, given $(\ell, h)$.
Follow the guidelines in Lecture Note 1 and the examples in Lecture Note 7.

For readability, use mathematical expressions (such as a sums, maximums, etc.)
that any competent programmer could implement in an obvious way,
without giving the details of how to compute those expressions.

\end{enumerate}

You can test your algorithm at \url{https://www.hackerrank.com/nealeyoung-algs101-schedule-and-valid-subsets}
(the ``Maximum-value schedule'' challenge).
(This is just for testing. Hacker-Rank submissions won't be graded.)

\continued

\setlength{\ITEMHT}{1.9in}
{\injectEnumerate
  \inputAnswer{problem_1_answer}
}

\newpage
%%%%%%%%%%%%%%%%%%%%%%%%%%%%%%%%%%%%%%%%%%%%%%%%%%%%%%%%%%%% 
%%%%%%%%%%%%%%%%%%%%%%%%%%%%%%%%%%%%%%%%%%%%%%%%%%%%%%%%%%%% 
%%%%%%%%%%%%%%%%%%%%%%%%%%%%%%%%%%%%%%%%%%%%%%%%%%%%%%%%%%%% 
%%%%%%%%%%%%%%%%%%%%%%%%%%%%%%%%%%%%%%%%%%%%%%%%%%%%%%%%%%%% 

%%%%% EDIT THE FILES NAMED problem_X_answer.tex.  DON'T EDIT THIS FILE 

%%%%%%%%%%%%%%%% PROBLEM 2 %%%%%%%%%%%%%%%%%

\problem\emph{(Lecture Note 7)}
Consider the following problem:

\begin{mylist}
\item[\bf input:] A universe $U$ (a set) and sequence $X = (x[1], x[2], \ldots, x[n]$), where each $x[i]$ is in $U$.
\item[\bf output:] The number of \emph{valid} subsets $S$ of $\{1,2,\ldots,n\}$,
  where $S$ is valid if the corresponding subsequence of $X$
  has no adjacent equal elements.  
  For example, if $U=\{a,b\}$ and $X=a,b,b,a$, then 11 of the 16 possible subsets are valid:
\end{mylist}
\centerline{
\begin{tabular}{rcll}
  subset  & subsequence & valid?  \\ \hline
 $\{1,2,3,4\}$ & $a, b, b, a$ & no \\
 $\{2,3,4\}$ & $-, b, b, a$ & no\\
 $\{1,3,4\}$ & $a, -, b, a$ & yes\\
 $\{1,2,4\}$ & $a, b, -, a$ & yes\\
 $\{3,4\}$ & $-, -, b, a$ & yes\\
 $\{2,4\}$ & $-, b, -, a$ & yes\\
 $\{1,4\}$ & $a, -, -, a$ & no\\
 $\{4\}$ & $-, -, -, a$ & yes\\ \hline
\end{tabular}~~~~{}
\begin{tabular}{rcll}
  subset  & subsequence & valid?  \\ \hline
 $\{1,2,3\}$ & $a, b, b,-$ & no \\
 $\{2,3\}$ & $-, b, b, -$ & no\\
 $\{1,3\}$ & $a, -, b, -$ & yes\\
 $\{3\}$ & $-, -, b, -$ & yes\\
 $\{1,2\}$ & $a, b, -, -$ & yes\\
 $\{2\}$ & $-, b, -, -$ & yes\\
 $\{1\}$ & $a, -, -, -$ & yes\\
 $\{\}$ & $-, -, -, -$ & yes \\\hline
\end{tabular}
}
\medskip\noindent
(If all elements of $X$ are the same, there are $n+1$ valid subsets.
If all are distinct, there are $2^n$.)

\paragraph{A dynamic-programming algorithm (with one small bug that you must find).}
Fix an input $X=x[1], \ldots, x[n]$.
Assume for notational convenience that $U$ contains some element that is not in $X$.
(This is without loss of generality; otherwise, just add an artificial element to $U$.)
Let $\star$ denote this element.

\paragraph{The subproblems to be solved.} 
For each $i\in\{0,1,\ldots,n\}$ and $u\in U$,
define $V(i, u)$ to be the collection of valid subsets of $\{1,2,\ldots,i\}$
such that the corresponding subsequence of $X$ ends with an element \emph{other than} $u$.
The desired output, for input $X$, is $|V(n, \star)|$.

\paragraph{The recurrence relation?}
\begin{lemma}
  For all $u\in U$ and each $i \in \{0, 1,\ldots,n\}$,
  \[
  |V(i, u)| =
  \begin{cases}
    1 & \text{if } i = 0, \\
    |V(i-1, u)| & \text{if } i>0 \text{ and } x[i] = u, \\
    |V(i-1, u)| + |V(i-1, \star)|& \text{if } i > 0 \text{ and } x[i] \ne u.
  \end{cases}
  \]
\end{lemma}
\begin{proof}[Proof sketch?]
  Consider any $u\in U$ and $i\in \{0, 1, \ldots, n\}$.
  For $i=0$, we have $V(0, u) = \{\emptyset\}$, so $|V(0, u)| = 1$. 
  So consider the remaining case, $i\ge 1$.
  
  In the subcase that $x[i] = u$,
  index $i$ is not in any subset in $V(i, u)$,
  so $V(i, u)$ is just $V(i-1, u)$, and the recurrence holds.

  In the remaining subcase, $x[i] \ne u$.
  Partition the subsets $S$ in $V(i, u)$ into two types:
  those that don't contain element $i$, and those that do.
  Those that don't are just the subsets in $V(i-1, u)$.
  Those that do are of the form $S\cup\{i\}$
  where $S$ is a subset in $V(i-1,\star)$.
  Thus, the total number of subsets in $V(i,u)$ is $|V(i-1, u)| + |V(i-1, \star)|$,
  which equals the right-hand side of the recurrence (in the subcase $x[i]\ne u$).
  \renewcommand{\qed}{\hfill \framebox{\mbox{\tiny ?}}}
\end{proof}

\continued

\paragraph{The algorithm.}
This gives the following iterative algorithm:

\newcommand{\countV}{\textsf{countV}}

\smallskip 

\begin{myframed}
  \begin{steps}
  \item[]
    def \countV$(U, X=(x[1], x[2], \ldots, x[n]))$:
    \comment{(assume $U$ contains some element~$\star \not\in X$)}
    \step for $i=0,1,2,\ldots, n$:
    \step~~~for $u\in U$:
      \step~~~~~~\(
      v[i, u] =
      \begin{cases}
        1 & \text{if } i=0, \\
        v[i-1, u] & \text{if } i > 0 \text{ and } x[i] = u, \\
        v[i-1, u] + v[i-1, \star]& \text{if } i > 0 \text{ and } x[i] \ne u.
      \end{cases}
      \)
    \step return $v[n, \star]$
  \end{steps}
\end{myframed}

\vspace*{-10pt}
\paragraph{Run time.}
By inspection, the running time is dominated by the time for the additions.
There are $O(n|U|)$ additions---at most one for each subproblem.
The integers that arise can be large, but no larger than $2^n$, so can be represented in $O(n)$ bits,
so each addition can be done in $O(n)$ time.
Hence, the algorithm runs in time $O(n^2|U|)$.



\noindent
\adjustbox{valign=t}{\parbox{4.75in}{%
    \paragraph{Your task: fix the bug!}
    Given the example input where $U=\{a,b,\star\}$ and $X=(a,b,b,a)$,
    the algorithm sets the array $v$ as shown to the right, then outputs 16.
    This is wrong! As discussed at the start, the correct output for this example is 11.
}}~~~{}
\begin{tabular}[t]{c|c@{~}c@{~}c@{~}c@{~}}
  $i$ & $x[i]$ & $v[i,a]$ & $v[i,b]$ & $v[i, \star]$ \\ \hline
  0    &          & 1 & 1 & 1  \\
  1    &  $a$  & 1 & 2 & 2 \\
  2    &  $b$  & 3 & 2 & 4 \\
  3    &  $b$  & 7 & 2 & 8 \\
  4    &  $a$  & 7 & 10 & 16
\end{tabular}

\begin{enumerate}[label=(\alph*)]

\item

  What is the smallest $i$ for which, on the example input, $v[i,\star]$
  as computed is not $|V(i, \star)|$?

\item
  
  What is the set $V(i,\star)$ for that $i$?

\item

  Give a corrected recurrence relation, making as few changes as possible.

  % \item

  %   Give a corrected proof sketch, making as few changes as possible.

\item

  Simulate the corrected algorithm on the example.
  Show the resulting array $v$.

\item

  Give pseudo-code for the correct algorithm, making as few changes as possible.

\end{enumerate}

You can test your algorithm at \url{https://www.hackerrank.com/nealeyoung-algs101-schedule-and-valid-subsets}
(the ``Number of valid subsets'' challenge, which asks for the number mod $2^{30}$).
(This is just for testing. Hacker-Rank submissions won't be graded.)

\newpage

\setlength{\ITEMHT}{1.9in}
{\injectEnumerate
  \inputAnswer{problem_2_answer}
}


%%%%%%%%%%%%%%%%%%%%%%%%%%%%%%%%%%%%%%%%%%%%%%%%%%%%%%%%%%%% 
%%%%%%%%%%%%%%%%%%%%%%%%%%%%%%%%%%%%%%%%%%%%%%%%%%%%%%%%%%%% 
%%%%%%%%%%%%%%%%%%%%%%%%%%%%%%%%%%%%%%%%%%%%%%%%%%%%%%%%%%%% 
%%%%%%%%%%%%%%%%%%%%%%%%%%%%%%%%%%%%%%%%%%%%%%%%%%%%%%%%%%%% 

%%%%% EDIT THE FILES NAMED problem_X_answer.tex.  DON'T EDIT THIS FILE 

%%%%%%%%%%%%%%%% PROBLEM 3 %%%%%%%%%%%%%%%%%

\newpage

\small
\problem\emph{(Lecture Note 7)}
Design a dynamic-programming algorithm for the following problem:

\noindent\textbf{input:} A sequence $(a_1, a_2, \ldots, a_n)$ of $n$ non-negative numbers, defining the following ``Game of Piles''.

  The game starts with a row of $n$ piles $1,2,\ldots,n$ of coins, where pile $i$
  contains coins of total value $a_i$.
  Alice and Bob take turns removing either the first or last of the remaining piles, and placing it in their personal stash, until there no piles remain.
  (If Alice moves pile $1$ into her stash, then Bob can move either pile $2$ or pile $n$ into his, and so on.)
  Alice's final score is the total value of her stash, minus the total value of Bob's.
  Bob's final score is the total value of his personal stash, minus the total value of Alice's.
  (For example, suppose $a=(10, 2, 8)$.  Alice can take the first pile to guarantee a score of at least $10 + 2 - 8 = 4$.
  But if she takes the last pile she could only guarantee a score of at least $8 + 2 - 10 = 0$.)

\noindent\textbf{output:}
The maximum score that the first player can guarantee for his or herself,
assuming the opponent plays optimally to maximize his score.
Each player's score is the negation of the opponents score, so by playing to maximize his or her own score, each player is also playing to minimize the opponent's score.

\begin{enumerate}[label=(\alph*)]
  \item

Per Section 7.4, define the subproblems your algorithm solves.
Which is the final answer?

\item

State an appropriate recurrence relation, including boundary cases.

\item
  
  Explain clearly in at most five lines why the recurrence relation is correct.

% Give a long-form proof of the lemma (that your recurrence relation is correct).

% \item

\item

Give the big-$\Theta$ run-time of the resulting algorithm (iterative, or recursive and memoized).

\item
  
Explain clearly in at most three lines why this bound holds.

\item
  
Give pseudo-code for an iterative implementation of the resulting algorithm.

\item

  Submit your algorithm to the Hacker-Rank challenge at the URL in the footnote below.\footnote
  {\url{https://www.hackerrank.com/nealeyoung-algs101-game-of-piles}}
  (For full credit on this part, your submission should pass all tests.)
  What is the username of your Hacker-Rank submission? 
  
\end{enumerate}

\paragraph{Clarification of what it means to ``play optimally''.} Consider the game with $a=(2, 2, 5, 3)$, with Alice as first player.
 
\begin{description}
\item[Option 1:]
  Suppose Alice starts by taking the first pile (value 2).

  She'll score 2 for that, and the game that remains will be 2, 5, 3, with Bob as first player.
  Considering just that game 2, 5, 3---whatever Bob does, Alice will take the 5, so Bob will end up with the 2 and the 3.  So Alice will score $5-2-3 = 0$ for this sub-game.
  So, if Alice starts the game $2, 2, 5, 3$ by taking the 2, her total score for the game will be $2 + 0 = 2$.
 
\item[Option 2:]
  Suppose Alice starts by taking the last pile (value 3).

  She'll score 3 for that, and the game that remains will be 2, 2, 5, with Bob as first player.
  Considering just that game $2, 2, 5$---Bob should take the 5, then Alice will take a 2, and Bob will take the other 2.  So Alice will score $2-5-2 = -5$ for this sub-game.
  So, if Alice starts the game $2, 2, 5, 3$ by taking the 3, her total score for the game will be $3 - 5 = -2$.
\end{description}

The optimal first move for Alice is to take the first pile (value 2).  This maximizes her score, assuming both players play the remaining game optimally.
This kind of reasoning defines ``optimal play'' for the general case.   Assuming we know what optimal play is for any game with, say, $n-1$ piles, we can inductively define what optimal play is for any game with $n$ piles.  Namely, the first player chooses the first move to maximize his or her resulting total score, assuming that after the first move, both players optimally play whatever game (on $n-1$ piles) remains.
 

\emph{Warning: you might be tempted to try a greedy algorithm that chooses the first move just by looking at just a few piles near each end.  This won't work.  You'll need to bring to bear the full power of dynamic programming!}

% Note that we say ``this maximizes her score\ldots assuming both players play the remaining game optimally''.  If Bob isn't too smart (e.g., if Bob always takes the first pile), then Alice may have a strategy that gives her an even better outcome (and makes Bob's worse).  But we aren't considering this case.  We are assuming that Bob always plays to maximize his score.

\normalsize

\newpage

\setlength{\ITEMHT}{1.9in}
{\injectEnumerate
  \inputAnswer{problem_3_answer}
}

%%%%%%%%%%%%%%%%%%%%%%%%%%%%%%%%%%%%%%%%%%%%%%%%%%%%%%%%%%%% 
%%%%%%%%%%%%%%%%%%%%%%%%%%%%%%%%%%%%%%%%%%%%%%%%%%%%%%%%%%%% 
%%%%%%%%%%%%%%%%%%%%%%%%%%%%%%%%%%%%%%%%%%%%%%%%%%%%%%%%%%%% 
%%%%%%%%%%%%%%%%%%%%%%%%%%%%%%%%%%%%%%%%%%%%%%%%%%%%%%%%%%%% 

%%%%% EDIT THE FILES NAMED problem_X_answer.tex.  DON'T EDIT THIS FILE 

%%%%%%%%%%%%%%%% PROBLEM 3 %%%%%%%%%%%%%%%%%

\newpage

\problem\emph{(Exercise 5.13)}
Here is a greedy \prob{Minimum Spanning Tree} algorithm:

\begin{myframed}
  \begin{steps}
  \item[{{\sf cycleBreaker}$(G=(V,E))$:}]
    \comment{--- $G$ is a connected, undirected, edge-weighted graph}
    \step initialize $E' = E$
    \step while the edges in $E'$ contain a cycle:
    \begin{block}
      \step choose any cycle $C$ in $E'$, and let $(u, w)$ be an edge of maximum weight in $C$
      \step delete $(u, w)$ from $E'$
    \end{block}
    \step return the spanning tree formed by the edges in $E'$
  \end{steps}
\end{myframed}

Complete the incomplete proof of correctness for this algorithm that is given in the template.
You may wish to review the proof of correctness of Prim's MST algorithm in Section 5.9.1.
We will use the following observation:

% [review 2026-09-27] fix: removed \setcounter{observation}{6}, which numbered this observation to match an old numbering of Lecture Note 5

\begin{observation}\label{exchange}
  Let $(u,w)$ be any edge in any spanning tree $T$ of an undirected graph $G=(V,E)$.
  Removing $(u,w)$ splits $T$ into two connected subtrees $T_u$ and $T_w$ (see below).
  Let $(x,y)$ be any edge in $E$ with $x\in T_u$ and $y\in T_w$.
  Then removing $(u, w)$ from $T$ and adding $(x, y)$ 
  gives a spanning tree $T'=\{(x, y)\}\cup T \setminus \{(u, w)\}$ of $G$.
\end{observation}

\tikzset{% >= latex,
  vertex/.style ={circle, draw=black, minimum width = 12pt, inner sep=0pt},
  subtree/.style={
    draw,dotted,shape border uses incircle, yshift=30pt, % fill=lightgray,
    isosceles triangle,shape border rotate=90, % yshift=-1.45cm
    minimum height=40pt,
    minimum width=60pt,
    isosceles triangle stretches,
  }
}

\begin{center}
\begin{tikzpicture}
  \node[vertex] (v) at (3, 3.5) {}; 
  \node[vertex] (c1) at (2, 3) {}; 
  \node[vertex] (c2) at (1.5, 2) {$u$}; 
  \node[vertex] (c3) at (2.2, 1) {$w$}; 
  \node[vertex] (wp) at (3.5, 1) {}; 
  \node[vertex] (v1) at (2.9, 4.5) {};
  \node[vertex] (v2) at (4, 3.2) {}; 
  \node[vertex] (c11) at (1.4, 3.8) {}; 
  \node[vertex] (c111) at (0.5, 4.3) {}; 
  \node[vertex] (c112) at (0.4, 3.5) {}; 
  \node[vertex] (c21) at (0.5, 2) {}; 
  \node[vertex] (c211) at (-0.5, 3) {}; 
  \node[vertex] (c212) at (-0.5, 1) {}; 
  \node[vertex] (c31) at (2, 0) {}; 
  \node[vertex] (w1) at (4.5, 1) {}; 
  \node[vertex] (w2) at (3.5, 0) {}; 
  \node at (5.5, 4.5) {\Large $T~\longrightarrow~~T'$}; 
  \node at (2.7, 0.5) {\large $T_w$};
  \node at (1, 3) {\large $T_u$};
  % (r) edge[treeEdge, dotted] (u)
  % \draw (v) -- (wp); 
  \draw (v) -- (c1) -- (c2); 
  \draw[ultra thick] (c2) -- (c3); 
  \draw (c3) -- (wp); 
  \draw (v1) -- (v) -- (v2); 
  \draw (c1) -- (c11); 
  \draw (c2) -- (c21); 
  \draw (c3) -- (c31); 
  \draw (c111) -- (c11) -- (c112); 
  \draw (c211) -- (c21) -- (c212); 
  \draw (w1) -- (wp) -- (w2); 
  \draw[dotted] (0,0) to [bend left=20] (5,2.5); 
\end{tikzpicture}
\begin{tikzpicture}
  \node[vertex] (v) at (3, 3.5) {$x$}; 
  \node[vertex] (c1) at (2, 3) {}; 
  \node[vertex] (c2) at (1.5, 2) {$u$}; 
  \node[vertex] (c3) at (2.2, 1) {$w$}; 
  \node[vertex] (wp) at (3.5, 1) {$y$}; 
  \node[vertex] (v1) at (2.9, 4.5) {};
  \node[vertex] (v2) at (4, 3.2) {}; 
  \node[vertex] (c11) at (1.4, 3.8) {}; 
  \node[vertex] (c111) at (0.5, 4.3) {}; 
  \node[vertex] (c112) at (0.4, 3.5) {}; 
  \node[vertex] (c21) at (0.5, 2) {}; 
  \node[vertex] (c211) at (-0.5, 3) {}; 
  \node[vertex] (c212) at (-0.5, 1) {}; 
  \node[vertex] (c31) at (2, 0) {}; 
  \node[vertex] (w1) at (4.5, 1) {}; 
  \node[vertex] (w2) at (3.5, 0) {}; 
  % (r) edge[treeEdge, dotted] (u)
  \draw (c1) -- (c2); 
  % \draw (c2) -- (c3);
  \draw (c3) -- (wp); 
  \draw[ultra thick] (wp) -- (v); 
  \draw (v) -- (c1); 
  \draw (v1) -- (v) -- (v2); 
  \draw (c1) -- (c11); 
  \draw (c2) -- (c21); 
  \draw (c3) -- (c31); 
  \draw (c111) -- (c11) -- (c112); 
  \draw (c211) -- (c21) -- (c212); 
  \draw (w1) -- (wp) -- (w2); 
  \draw[dotted] (0,0) to [bend left=20] (5,2.5); 
\end{tikzpicture}
\end{center}


\newpage

\setlength{\ITEMHT}{1.9in}
{\injectEnumerate
  \inputAnswer{problem_4_answer}
}

\end{document}

%%% Local Variables:
%%% mode: latex
%%% TeX-master: t
%%% End:
